\documentclass[10pt,a4paper]{amsart}
\usepackage[margin=19mm]{geometry}
\usepackage{amsmath,amssymb,mathtools}
\usepackage[scr=rsfs,cal=euler]{mathalfa}
\usepackage{xfrac}
\usepackage[unicode,hidelinks,bookmarks=false]{hyperref}
\newcommand{\ie}{i.e.\ }
\newcommand{\cf}{cf.\ }
\newcommand{\resp}{respectively\ }
\newcommand{\Subsection}{\subsection}
\providecommand{\nobreakdash}{\nobreak\hbox{-}}
\setlength{\emergencystretch}{2em}
\multlinegap0pt
\usepackage{amssymb}
\usepackage{tikz}
\newtheorem{theorem}{Theorem}
\newtheorem{lemma}{Lemma}
\usepackage{cleveref}
\crefname{theorem}{Theorem}{Theorems}
\crefname{lemma}{Lemma}{Lemmas}
\newcommand{\II}[1]{\mathbf{1}{#1}}
\newcommand{\abs}[1]{\left\lvert#1\right\rvert}
\newcommand{\dd}{\mathrm{d}}
\title{Uniform Geodesic Drawings of Graphs}
\author{Saba Lepsveridze, Oriol Sol\'e-Pi}
\date{Source: arXiv:2605.16700v1 (CC BY 4.0)}
\begin{document}
\maketitle

\noindent\emph{Source and licence.} Uniform Geodesic Drawings of Graphs. Version arXiv:2605.16700v1 (CC BY 4.0), distributed by arXiv under the Creative Commons Attribution 4.0 International licence, http://creativecommons.org/licenses/by/4.0/, as recorded in the arXiv licence metadata for this exact version and shown on the abstract page https://arxiv.org/abs/2605.16700v1. Full licence notice: \texttt{licenses/arxiv-2605.16700-CC-BY-4.0.txt}.\par\smallskip

\noindent\textit{Editorial scope.} Counted unit: the article's theorem thm:continuous-spherical-crossing-main (source0.tex lines 213--220). The article's complete original proof of it is delivered with it (source0.tex lines 335--400, one \texttt{proof} environment of 66 lines, which the article places away from the statement). No mathematics is written, repaired or re-derived: the statement and the proof are the article's own lines, in the article's own notation, and no retained line is substituted or re-flowed.  Retained uncounted as the source-local context the proof needs: lines 293--298; lines 300--302; lines 304--315; lines 321--328.  Nothing was omitted from any retained span, so the package carries no omission record.  No retained line contains a citation, so the wrapper carries no bibliography and no bibliography entry.  No retained line contains a figure environment, an image-inclusion command or drawing code, so no image is delivered and the package declares no asset dependency.  Editorial formatting only, in addition: an AMS \texttt{amsart} preamble that restates the article's own declarations and macros used by the retained lines, namely the environments \texttt{theorem}, \texttt{lemma} on separate counters, together with the standard packages those lines need.  The retained environments print in this document as Theorem 1, Lemma 1 in order of appearance, whereas the article prints the same items with the numbering of its own full document.  The complete native source of the article as distributed in the arXiv e-print for this exact version is bundled unchanged as \texttt{sources/200740/source0.tex}.
\begin{equation}\label{eq:normalized-flux-crossing-representation}
        \operatorname{Cr}(w)=
        \int_{{\mathbb S}^2}\int_0^{2\pi}\int_0^{2\pi}
        g(x,\theta)g(x,\phi)\abs{\sin(\theta-\phi)}
        \,\dd\theta\,\dd\phi\,\sigma(\dd x)\,.
\end{equation}

\begin{equation}\label{eq:normalized-incidence-identity}
       e(w)= \int_{{\mathbb S}^2}\int_0^{2\pi}a(x,\theta)\,\dd\theta\,\sigma(\dd x)\,.
\end{equation}

\begin{lemma}\label{lem:normalized-bathtub-main}
For almost every \((x,\theta) \in {\mathbb S}^2 \times [0, 2\pi)\), we have 
\begin{equation}\label{eq:normalized-pointwise-bathtub}
    g(x,\theta) \geq \varphi (a (x,\theta))\,,
\end{equation}
where $\varphi : [0,1/(8\pi^2)] \to {\mathbb R}$ is defined by 
\begin{equation}\label{eq:normalized-bathtub-profile}
        \varphi(\alpha):=
        \frac{\sin \tau(\alpha)-\tau(\alpha)\cos \tau(\alpha)}{(4\pi)^2}\,,
\end{equation}
where $\tau(\alpha) \in [0,\pi]$ is determined by $16\pi^2 \alpha =1-\cos  \tau(\alpha).$
\end{lemma}

\begin{equation}\label{eq:busemann-circle-main}
        \int_0^{2\pi}\int_0^{2\pi}
        f(\theta)f(\phi)\abs{\sin(\theta-\phi)}
        \,\dd\theta\,\dd\phi
        \ge
        \frac1{\pi^2}
        \left(\int_0^{2\pi}f(\theta)^{2/3}\,\dd\theta\right)^3.
\end{equation}

\begin{theorem}\label{thm:continuous-spherical-crossing-main}
Let \(w:{\mathbb S}^2\times{\mathbb S}^2\to[0,1]\) be measurable and symmetric. Then,
\begin{equation}\label{eq:continuous-spherical-bound-main}
        \operatorname{Cr}(w)\ge
        \frac{(\sin t-t\cos t)^2}{8\pi^2}\,,
\end{equation}
where \(t\in[0,\pi]\) satisfies $\cos t = 1-2e(w)$. Equality is attained by \( w(x,y):=\II{\{d(x,y)\le t\}}.\) 
\end{theorem}

\begin{proof}[Proof of \Cref{thm:continuous-spherical-crossing-main}]
The crossing representation \eqref{eq:normalized-flux-crossing-representation} and Busemann's inequality \eqref{eq:busemann-circle-main} give
\begin{align*}
    \operatorname{Cr}(w)&=\int_{{\mathbb S}^2}\int_0^{2\pi}\int_0^{2\pi}
        g(x,\theta)g(x,\phi)\abs{\sin(\theta-\phi)}
        \,\dd\theta\,\dd\phi\,\sigma(\dd x) \\
        &\geq  \frac1{\pi^2}\int_{{\mathbb S}^2} 
        \left(\int_0^{2\pi}g(x,\theta)^{2/3}\,\dd\theta\right)^3 \sigma(\dd x)\,.
\end{align*}
H\"older's inequality on \(({\mathbb S}^2,\sigma)\) with \(\sigma({\mathbb S}^2)=4\pi\) yields
\begin{align*}
    \operatorname{Cr}(w)&\geq  \frac1{\pi^2}\int_{{\mathbb S}^2} \left(\int_0^{2\pi}g(x,\theta)^{2/3}\,\dd\theta\right)^3 \sigma(\dd x) \geq \frac1{16\pi^4}\left(\int_{{\mathbb S}^2} \int_0^{2\pi}g(x,\theta)^{2/3}\,\dd\theta \,\sigma(\dd x) \right)^3.
\end{align*}
Let $h(\alpha) = \varphi(\alpha)^{2/3}$. Then, the pointwise estimate \Cref{lem:normalized-bathtub-main} gives 
\[
        g(x,\theta)^{2/3}\ge h(a(x,\theta))\,.
\]
By the convexity of \(h\), Jensen's inequality on the space
\({\mathbb S}^2\times[0,2\pi)\) implies
\begin{equation}\label{eq:main-proof-jensen-step}   \int_{{\mathbb S}^2}\int_0^{2\pi}g(x,\theta)^{2/3}\,\dd\theta\,\sigma(\dd x)
        \ge
        8\pi^2\,
        h\left(
        \frac1{8\pi^2}
        \int_{{\mathbb S}^2}\int_0^{2\pi}a(x,\theta)\,\dd\theta\,\sigma(\dd x)
        \right)  =8\pi^2\,h\left(\frac{e(w)}{8\pi^2}\right)\,,
\end{equation}
where the last step uses the identity \eqref{eq:normalized-incidence-identity}.
Let \(t\in[0,\pi]\) be determined by \(\cos t=1-2e(w)\).  The definition of \(\varphi\) gives
\[
        h\left(\frac{e(w)}{8\pi^2}\right)
        =
        \left(\frac{\sin t-t\cos t}{(4\pi)^2}\right)^{2/3}\,.
\]
Substituting this into the preceding bound and using \eqref{eq:main-proof-jensen-step}, we obtain
\[
        \operatorname{Cr}(w)
        \ge
        \frac1{16\pi^4}
        (8\pi^2)^3
        \left(\frac{\sin t-t\cos t}{(4\pi)^2}\right)^2
        =
        \frac{(\sin t-t\cos t)^2}{8\pi^2}\,.
\]
This proves the lower bound. It remains to check sharpness.  For density \(w(x,y)=\II{\{d(x,y)\le t\}}\), we have \(e(w)=(1-\cos t)/2\).  Moreover, we have that
\[
        g(x,\theta)
        =
        \frac1{(4\pi)^2}\int_0^t u\sin u\,\dd u
        =
        \frac{\sin t-t\cos t}{(4\pi)^2}
\]
independently of \((x,\theta)\).  Since $\sigma({\mathbb S}^2)=4\pi$ and 
\[
        \int_0^{2\pi}\int_0^{2\pi}\abs{\sin(\theta-\phi)}\,\dd\theta\,\dd\phi=8\pi\,,
\]
the flux formula \eqref{eq:normalized-flux-crossing-representation} gives
\[
        \operatorname{Cr}(w)
        =32\pi^2
        \left(\frac{\sin t-t\cos t}{(4\pi)^2}\right)^2
        =
        \frac{(\sin t-t\cos t)^2}{8\pi^2}\,.
\]
Hence equality is attained by \(w\).
\end{proof}

\end{document}
